\section{Executive Summary}
\label{sec:summary}

\paragraph{Trading concept.} Equity markets price AI mostly through the firms that sell it; the firms and industries that \emph{use} it are priced on a promise of future productivity. Our thesis is that the macro data already tell us which industries are using AI to replace workers and which are using it to make workers more productive, and that this distinction matters for cash flows. Where AI substitutes for labor, unit labor costs fall and margins rise; where it complements labor, the gain is bid away in wages. We ask whether the stock market separates the two.

\paragraph{Macro events.} The information arrives through four scheduled releases: BEA's quarterly and annual GDP-by-industry tables (output, compensation, and prices by industry), the monthly BLS Employment Situation (employment and hours by industry), Census's Quarterly Workforce Indicators (employment and hiring by age, published with roughly a one-year lag), and the monthly CPS unemployment rates by age. None of these releases moves markets on the day the way CPI or payrolls do; that is part of the point.

\paragraph{Target assets.} Sixteen U.S. sector portfolios built from the Fama--French 49 industries, traded long-short; the tradable equivalent is a set of sector ETFs.

\paragraph{Core strategy.} Rank sectors each month on a substitution score (how fast juniors are losing ground relative to prime-age workers in that sector), go long the top ranks and short the bottom ranks with dollar-neutral rank weights, and size the position with a real-time Bayesian posterior that starts agnostic about the sign. Interest-rate exposure is measured explicitly, because AI-exposed sectors behave like long-duration assets.

\paragraph{What we find.} The real-side evidence is strong and the pricing evidence is not. AI-exposed sectors show (i) no acceleration in measured output per hour, (ii) a relative decline of roughly 7 log points in junior employment per standard deviation of exposure, with an additional post-ChatGPT decline of about 1.2 log points a year, and (iii) a falling labor share, about 0.8 percentage points per year per standard deviation of exposure (bootstrap $p<0.01$). None of these patterns appears in the IT boom. The long-short strategy, however, earns a Sharpe ratio of $-0.19$ over 1996--2026 and no significant alpha. Our recommendation is to treat the substitution signal as a risk lens on AI-exposed equities rather than as a stand-alone alpha source, for reasons laid out in Sections~\ref{sec:forward} and~\ref{sec:strategy}.
